%----------------------------------------------------------%
%              VERIFICA DELLA CIRCONFERENZA                %
%----------------------------------------------------------%
\begin{question}[ID=1.1, subtitle={Verifica della circonferenza $x^2+y^2-4x+6y-3=0$}, difficulty=facile, tags={circonferenza}]
Verifica se l'equazione $x^2+y^2-4x+6y-3=0$ rappresenta una circonferenza e determina eventualmente centro e raggio.
\begin{solution}$C(2,-3),\ r=4$\end{solution}
\end{question}

%----------------------------------------------------------%

\begin{question}[ID=1.2, subtitle={Verifica di una circonferenza non reale}, difficulty=facile, tags={circonferenza}]
Verifica se l'equazione $x^2+y^2+2x-8y+20=0$ rappresenta una circonferenza.
\begin{solution}L'equazione non rappresenta una circonferenza reale.\end{solution}
\end{question}

%----------------------------------------------------------%

\begin{question}[ID=1.3, subtitle={Valori di $k$ per la forma canonica della circonferenza}, difficulty=media, tags={circonferenza, parametri}]
Studia per quali valori di $k$ l'equazione $x^2+y^2+(k-1)x+2y-5=0$ rappresenta una circonferenza.
\begin{solution}$\forall k\in\mathbb{R}$\end{solution}
\end{question}

%----------------------------------------------------------%

\begin{question}[ID=1.4, subtitle={Intervallo di $k$ per cui esiste la circonferenza}, difficulty=media, tags={circonferenza, parametri}]
Studia per quali valori di $k$ l'equazione $x^2+y^2+2(k-2)x+(k+1)y+1=0$ rappresenta una circonferenza.
\begin{solution}$k<\dfrac{7-2\sqrt{11}}{5} \vee k>\dfrac{7+2\sqrt{11}}{5}$\end{solution}
\end{question}

%----------------------------------------------------------%

\begin{question}[ID=1.5, subtitle={Raggio costante al variare di $k$}, difficulty=media, tags={circonferenza, parametri}]
Studia per quali valori di $k$ l'equazione $x^2+y^2-2kx+4y+k^2-9=0$ rappresenta una circonferenza.
\begin{solution}$\forall k\in\mathbb{R},\ r=\sqrt{13}$\end{solution}
\end{question}

%----------------------------------------------------------%
%               RICERCA DELLA CIRCONFERENZA                %
%----------------------------------------------------------%
\begin{question}[ID=2.1, subtitle={Circonferenza con centro nell'origine e $r=3$}, difficulty=facile, tags={circonferenza}]
Scrivi l'equazione della circonferenza con centro $C(0,0)$ e raggio $r=3$.
\begin{solution}$x^2+y^2=9$\end{solution}
\end{question}

%----------------------------------------------------------%

\begin{question}[ID=2.2, subtitle={Circonferenza con centro $C(2,-1)$ e $r=4$}, difficulty=facile, tags={circonferenza}]
Scrivi l'equazione della circonferenza con centro $C(2,-1)$ e raggio $r=4$.
\begin{solution}$(x-2)^2+(y+1)^2=16$\end{solution}
\end{question}

%----------------------------------------------------------%

\begin{question}[ID=2.3, subtitle={Verifica dell'appartenenza del punto $P(3,0)$}, difficulty=facile, tags={circonferenza}]
Verifica se il punto $P(3,0)$ appartiene alla circonferenza $x^2 + y^2 = 9$.
\begin{solution}Sì, $P(3,0)$ appartiene alla circonferenza.\end{solution}
\end{question}

%----------------------------------------------------------%

\begin{question}[ID=2.4, subtitle={Centro e raggio di $x^2+y^2-4x+6y=3$}, difficulty=facile, tags={circonferenza}]
Trova il centro e il raggio della circonferenza $x^2 + y^2 - 4x + 6y = 3$.
\begin{solution}$C(2,-3),\ r=4$\end{solution}
\end{question}

%----------------------------------------------------------%

\begin{question}[ID=2.5, subtitle={Circonferenza avente un segmento come diametro}, difficulty=media, tags={circonferenza}]
Scrivi l'equazione della circonferenza avente diametro con estremi $A(1,2)$ e $B(5,6)$.
\begin{solution}$(x-3)^2+(y-4)^2=8$\end{solution}
\end{question}

%----------------------------------------------------------%

\begin{question}[ID=2.6, subtitle={Circonferenza per tre punti $A(0,0)$, $B(2,0)$, $C(0,2)$}, difficulty=media, tags={circonferenza}]
Determina l'equazione della circonferenza passante per i punti $A(0,0)$, $B(2,0)$, $C(0,2)$.
\begin{solution}$(x-1)^2+(y-1)^2=2$\end{solution}
\end{question}

%----------------------------------------------------------%

\begin{question}[ID=2.7, subtitle={Circonferenza con centro sull'asse $x$ per $A(1,2)$ e $B(3,4)$}, difficulty=media, tags={circonferenza}]
Scrivi l'equazione della circonferenza con centro sull'asse $x$, passante per $A(1,2)$ e $B(3,4)$.
\begin{solution}$(x-2)^2+y^2=5$\end{solution}
\end{question}

%----------------------------------------------------------%

\begin{question}[ID=2.8, subtitle={Circonferenza per $A(1,0)$ e $B(0,1)$ con centro su $x=y$}, difficulty=media, tags={circonferenza}]
Trova la circonferenza passante per $A(1,0)$, $B(0,1)$ con centro sulla retta $x=y$.
\begin{solution}$\left(x-\tfrac{1}{2}\right)^2+\left(y-\tfrac{1}{2}\right)^2=\tfrac{1}{2}$\end{solution}
\end{question}

%----------------------------------------------------------%

\begin{question}[ID=2.9, subtitle={Circonferenza per $A(2,0)$, $B(0,2)$, $C(2,2)$}, difficulty=media, tags={circonferenza}]
Determina la circonferenza passante per $A(2,0)$, $B(0,2)$, $C(2,2)$.
\begin{solution}$(x-1)^2+(y-1)^2=1$\end{solution}
\end{question}

%----------------------------------------------------------%

\begin{question}[ID=2.10, subtitle={Circonferenza con centro sull'asse $y$ tangente all'asse $x$}, difficulty=media, tags={circonferenza}]
Trova la circonferenza con centro sull'asse $y$, passante per $A(2,1)$ e tangente all'asse $x$.
\begin{solution}$x^2+(y-1)^2=4$\end{solution}
\end{question}

%----------------------------------------------------------%

\begin{question}[ID=2.11, subtitle={Circonferenza per $A(0,0)$ e $B(4,0)$ tangente all'asse $y$}, difficulty=media, tags={circonferenza}]
Scrivi la circonferenza passante per $A(0,0)$, $B(4,0)$ e tangente all'asse $y$.
\begin{solution}$(x-2)^2+(y-2)^2=8$\end{solution}
\end{question}

%----------------------------------------------------------%

\begin{question}[ID=2.12, subtitle={Circonferenza per $A(1,1)$, $B(3,1)$, $C(2,4)$}, difficulty=difficile, tags={circonferenza}]
Determina la circonferenza passante per $A(1,1)$, $B(3,1)$, $C(2,4)$.
\begin{solution}$(x-2)^2+(y-2)^2=2$\end{solution}
\end{question}

%----------------------------------------------------------%

\begin{question}[ID=2.13, subtitle={Circonferenza con centro su $y=x$ per $A(2,0)$ e $B(0,2)$}, difficulty=difficile, tags={circonferenza}]
Trova la circonferenza con centro sulla retta $y=x$ passante per $A(2,0)$ e $B(0,2)$.
\begin{solution}$(x-1)^2+(y-1)^2=2$\end{solution}
\end{question}

%----------------------------------------------------------%

\begin{question}[ID=2.14, subtitle={Circonferenza per $A(0,0)$, $B(2,4)$, $C(4,0)$}, difficulty=difficile, tags={circonferenza}]
Determina la circonferenza passante per $A(0,0)$, $B(2,4)$, $C(4,0)$.
\begin{solution}$(x-2)^2+(y-1)^2=5$\end{solution}
\end{question}

%----------------------------------------------------------%

\begin{question}[ID=2.15, subtitle={Circonferenza per $A(1,0)$, $B(3,2)$, $C(5,0)$}, difficulty=difficile, tags={circonferenza}]
Trova la circonferenza passante per $A(1,0)$, $B(3,2)$, $C(5,0)$.
\begin{solution}$(x-3)^2+(y-1)^2=2$\end{solution}
\end{question}

%----------------------------------------------------------%

\begin{question}[ID=2.16, subtitle={Circonferenza per $A(0,1)$, $B(2,3)$, $C(4,1)$}, difficulty=difficile, tags={circonferenza}]
Determina la circonferenza passante per $A(0,1)$, $B(2,3)$, $C(4,1)$.
\begin{solution}$(x-2)^2+(y-2)^2=2$\end{solution}
\end{question}

%----------------------------------------------------------%
%      POSIZIONE RECIPROCA TRA RETTA E CIRCONFERENZA       %
%----------------------------------------------------------%
\begin{question}[ID=3.1, subtitle={Posizione di $y=2$ rispetto a $x^2+y^2=9$}, difficulty=facile, tags={retta, circonferenza, posizione reciproca}]
Determina la posizione reciproca tra la retta $y=2$ e la circonferenza $x^2+y^2=9$.
\begin{solution}Secante\end{solution}
\end{question}

%----------------------------------------------------------%

\begin{question}[ID=3.2, subtitle={Posizione di $x=3$ rispetto a $x^2+y^2=9$}, difficulty=facile, tags={retta, circonferenza, posizione reciproca}]
Determina la posizione reciproca tra la retta $x=3$ e la circonferenza $x^2+y^2=9$.
\begin{solution}Tangente\end{solution}
\end{question}

%----------------------------------------------------------%

\begin{question}[ID=3.3, subtitle={Posizione di $y=x$ rispetto a $x^2+y^2=4$}, difficulty=facile, tags={retta, circonferenza, posizione reciproca}]
Stabilisci la posizione tra la retta $y=x$ e la circonferenza $x^2+y^2=4$.
\begin{solution}Secante\end{solution}
\end{question}

%----------------------------------------------------------%

\begin{question}[ID=3.4, subtitle={Intersezioni di $y=mx$ con la circonferenza unitaria al variare di $m$}, difficulty=media, tags={retta, circonferenza, parametrico}]
La retta $y=mx$ interseca la circonferenza $x^2+y^2=1$. Studia il numero di intersezioni al variare di $m$.
\begin{solution}Sempre secante per ogni $m\in\mathbb{R}$\end{solution}
\end{question}

%----------------------------------------------------------%

\begin{question}[ID=3.5, subtitle={Posizione di $y=k$ rispetto a $x^2+y^2=4$ al variare di $k$}, difficulty=media, tags={parametrico, retta, circonferenza}]
Studia la posizione reciproca tra la retta $y=k$ e la circonferenza $x^2+y^2=4$ al variare di $k$.
\begin{solution}$|k|<2$ secante; $|k|=2$ tangente; $|k|>2$ esterna\end{solution}
\end{question}

%----------------------------------------------------------%

\begin{question}[ID=3.6, subtitle={Valori di $k$ per cui $y=x+k$ è tangente a $x^2+y^2=1$}, difficulty=media, tags={parametrico, tangente, circonferenza}]
Per quali valori di $k$ la retta $y=x+k$ è tangente alla circonferenza $x^2+y^2=1$?
\begin{solution}$k=\pm\sqrt{2}$\end{solution}
\end{question}

%----------------------------------------------------------%

\begin{question}[ID=3.7, subtitle={Posizione di $y=mx+1$ rispetto a $x^2+y^2=4$ al variare di $m$}, difficulty=facile, tags={parametrico, retta, circonferenza}]
Studia la posizione tra la retta $y=mx+1$ e la circonferenza $x^2+y^2=4$ al variare di $m$.
\begin{solution}Sempre secante per ogni $m\in\mathbb{R}$\end{solution}
\end{question}

%----------------------------------------------------------%

\begin{question}[ID=3.8, subtitle={Valori di $k$ per cui $y=2x+k$ è esterna a $x^2+y^2=1$}, difficulty=media, tags={parametrico, retta, circonferenza}]
Per quali valori di $k$ la retta $y=2x+k$ è esterna alla circonferenza $x^2+y^2=1$?
\begin{solution}$|k|>\sqrt{5}$\end{solution}
\end{question}

%----------------------------------------------------------%

\begin{question}[ID=3.9, subtitle={Valori di $k$ per cui $kx+y=1$ è tangente alla circonferenza unitaria}, difficulty=difficile, tags={parametrico, tangente}]
Determina per quali valori di $k$ la retta $kx+y=1$ è tangente alla circonferenza $x^2+y^2=1$.
\begin{solution}$k=0$\end{solution}
\end{question}

%----------------------------------------------------------%

\begin{question}[ID=3.10, subtitle={Posizione di $y=kx$ rispetto a $x^2+y^2=1$ e tangenza}, difficulty=difficile, tags={parametrico, retta, circonferenza}]
Studia la posizione reciproca tra la retta $y=kx$ e la circonferenza $x^2+y^2=1$ e determina per quali $k$ è tangente.
\begin{solution}Nessun valore di $k$ rende la retta tangente\end{solution}
\end{question}

%----------------------------------------------------------%

\begin{question}[ID=3.11, subtitle={Corda di lunghezza 4 con retta orizzontale sulla circonferenza}, difficulty=media, tags={circonferenza, retta, corda}]
Determina le equazioni delle rette parallele all'asse $x$ che individuano sulla circonferenza $x^2+y^2-4x-2y-4=0$ una corda di lunghezza $4$.
\begin{solution}$y=1\pm\sqrt{3}$\end{solution}
\end{question}

%----------------------------------------------------------%

\begin{question}[ID=3.12, subtitle={Corda di lunghezza $2\sqrt{2}$ parallela alla bisettrice}, difficulty=media, tags={circonferenza, corda}]
Determina le equazioni delle rette parallele alla bisettrice del primo e terzo quadrante che individuano sulla circonferenza $x^2+y^2-2x-2y-3=0$ una corda di lunghezza $2\sqrt{2}$.
\begin{solution}$y=x+1,\ y=x-3$\end{solution}
\end{question}

%----------------------------------------------------------%

\begin{question}[ID=3.13, subtitle={Corda di lunghezza 6 parallela a $y=2x$ sulla circonferenza}, difficulty=media, tags={circonferenza, corda}]
Determina le equazioni delle rette parallele alla retta $y=2x$ che individuano sulla circonferenza $x^2+y^2-6x+2y-3=0$ una corda di lunghezza $6$.
\begin{solution}$y=2x+3,\ y=2x-9$\end{solution}
\end{question}

%----------------------------------------------------------%

\begin{question}[ID=3.14, subtitle={Corda di lunghezza $2\sqrt{5}$ parallela a $x-y+1=0$}, difficulty=media, tags={circonferenza, corda}]
Determina le equazioni delle rette parallele alla retta $x-y+1=0$ che individuano sulla circonferenza $x^2+y^2-2x+4y-4=0$ una corda di lunghezza $2\sqrt{5}$.
\begin{solution}$x-y+1\pm\sqrt{5}=0$\end{solution}
\end{question}

%----------------------------------------------------------%

\begin{question}[ID=3.15, subtitle={Corda di lunghezza 4 parallela a $y=-x$ sulla circonferenza}, difficulty=difficile, tags={circonferenza, corda, parametrico}]
Determina le equazioni delle rette parallele alla retta $y=-x$ che individuano sulla circonferenza $x^2+y^2-4x+6y-3=0$ una corda di lunghezza $4$.
\begin{solution}$y=-x-1\pm 2\sqrt{6}$\end{solution}
\end{question}

%----------------------------------------------------------%
%          RETTA TANGENTE PASSANTE PER UN PUNTO            %
%----------------------------------------------------------%
\begin{question}[ID=4.1, subtitle={Tangente a $x^2+y^2=9$ nel punto $P(3,0)$}, difficulty=facile, tags={tangente, circonferenza, punto}]
Scrivi l'equazione della retta tangente alla circonferenza $x^2+y^2=9$ nel punto $P(3,0)$.
\begin{solution}$x=3$\end{solution}
\end{question}

%----------------------------------------------------------%

\begin{question}[ID=4.2, subtitle={Tangente a $x^2+y^2=4$ nel punto $P(0,2)$}, difficulty=facile, tags={tangente, circonferenza}]
Scrivi la retta tangente alla circonferenza $x^2+y^2=4$ nel punto $P(0,2)$.
\begin{solution}$y=2$\end{solution}
\end{question}

%----------------------------------------------------------%

\begin{question}[ID=4.3, subtitle={Tangente a $x^2+y^2=25$ nel punto $P(3,4)$}, difficulty=media, tags={tangente, circonferenza}]
Trova l'equazione della tangente alla circonferenza $x^2+y^2=25$ nel punto $P(3,4)$.
\begin{solution}$3x+4y=25$\end{solution}
\end{question}

%----------------------------------------------------------%

\begin{question}[ID=4.4, subtitle={Tangente a $x^2+y^2=10$ nel punto $P(1,3)$}, difficulty=media, tags={tangente, circonferenza}]
Determina la retta tangente alla circonferenza $x^2+y^2=10$ nel punto $P(1,3)$.
\begin{solution}$x+3y=10$\end{solution}
\end{question}

%----------------------------------------------------------%

\begin{question}[ID=4.5, subtitle={Tangente a $x^2+y^2=13$ nel punto $P(2,3)$}, difficulty=media, tags={tangente, circonferenza}]
Trova l'equazione della tangente alla circonferenza $x^2+y^2=13$ nel punto $P(2,3)$ e verifica che il punto appartenga alla circonferenza.
\begin{solution}$2x+3y=13$\end{solution}
\end{question}

%----------------------------------------------------------%
%         POSIZIONE RECIPROCA TRA 2 CIRCONFERENZE          %
%----------------------------------------------------------%
\begin{question}[ID=5.1, subtitle={Posizione di $C_1: x^2+y^2=9$ e $C_2: x^2+y^2=4$}, difficulty=facile, tags={circonferenze, posizione reciproca}]
Determina la posizione reciproca tra le circonferenze:
$C_1: x^2+y^2=9$ e $C_2: x^2+y^2=4$.
\begin{solution}Circonferenze concentriche\end{solution}
\end{question}

%----------------------------------------------------------%

\begin{question}[ID=5.2, subtitle={Posizione di $C_1: x^2+y^2=9$ e $C_2: (x-5)^2+y^2=4$}, difficulty=facile, tags={circonferenze, posizione reciproca}]
Determina la posizione tra:
$C_1: x^2+y^2=9$ e $C_2: (x-5)^2+y^2=4$.
\begin{solution}Tangenti esternamente\end{solution}
\end{question}

%----------------------------------------------------------%

\begin{question}[ID=5.3, subtitle={Posizione di $C_1: (x-1)^2+y^2=4$ e $C_2: (x-6)^2+y^2=9$}, difficulty=media, tags={circonferenze, posizione reciproca}]
Studia la posizione tra:
$C_1: (x-1)^2+y^2=4$ e $C_2: (x-6)^2+y^2=9$.
\begin{solution}Tangenti esternamente\end{solution}
\end{question}

%----------------------------------------------------------%

\begin{question}[ID=5.4, subtitle={Posizione di $C_1: x^2+y^2=25$ e $C_2: (x-3)^2+y^2=9$}, difficulty=media, tags={circonferenze, posizione reciproca}]
Determina la posizione tra:
$C_1: x^2+y^2=25$ e $C_2: (x-3)^2+y^2=9$.
\begin{solution}Secanti\end{solution}
\end{question}

%----------------------------------------------------------%

\begin{question}[ID=5.5, subtitle={Posizione reciproca di due circonferenze al variare di $k$}, difficulty=media, tags={circonferenze, posizione reciproca, parametrico}]
Studia la posizione reciproca al variare di $k$ tra:
$C_1: x^2+y^2=4$ e $C_2: (x-k)^2+y^2=1$.
\begin{solution}Secanti: $1<|k|<3$; tangenti interne: $|k|=1$; tangenti esterne: $|k|=3$; interne: $|k|<1$; esterne: $|k|>3$\end{solution}
\end{question}

%----------------------------------------------------------%
%                  FASCI DI CIRCONFERENZE                  %
%----------------------------------------------------------%
\begin{question}[ID=6.1, subtitle={Fascio $x^2+y^2+2x-4y+k=0$: centro, raggio e $k$ per l'origine}, difficulty=facile, tags={fasci, circonferenze}]
Considera il fascio $x^2+y^2+2x-4y+k=0$.
\begin{enumerate}
\item Determina il centro e il raggio della circonferenza.
\item Trova il valore di $k$ per cui la circonferenza passa per l'origine.
\end{enumerate}
\begin{solution}$C(-1,2),\ r=\sqrt{5-k},\ k=0$\end{solution}
\end{question}

%----------------------------------------------------------%

\begin{question}[ID=6.2, subtitle={Fascio $x^2+y^2-4x+k(x^2+y^2-4y)=0$: centro e degenere}, difficulty=media, tags={fasci, circonferenze}]
Considera il fascio $x^2+y^2-4x+k(x^2+y^2-4y)=0$.
\begin{enumerate}
\item Determina il centro del fascio.
\item Trova $k$ per cui la circonferenza è tangente all'asse $x$.
\item Studia per quali $k$ la circonferenza è degenere.
\end{enumerate}
\begin{solution}$C\!\left(2,\dfrac{2k}{1+k}\right)$; tangente per $k=0$ o $k=-2$; degenere per $k=-1$\end{solution}
\end{question}

%----------------------------------------------------------%

\begin{question}[ID=6.3, subtitle={Fascio $(x^2+y^2-2x)+k(x^2+y^2+2y)=0$: punti base e asse radicale}, difficulty=media, tags={fasci, circonferenze}]
Considera il fascio $(x^2+y^2-2x) + k(x^2+y^2+2y)=0$.
\begin{enumerate}
\item Determina i punti base del fascio.
\item Trova l'asse radicale.
\item Determina per quali $k$ la circonferenza passa per $P(1,-1)$.
\end{enumerate}
\begin{solution}$C_1(1,0,r=1)$, $C_2(0,-1,r=1)$; asse radicale: $x+y=0$; $k=0$\end{solution}
\end{question}

%----------------------------------------------------------%

\begin{question}[ID=6.4, subtitle={Fascio $x^2+y^2-6x+k(x^2+y^2-2y)=0$: centro e tangenza}, difficulty=media, tags={fasci, circonferenze}]
Considera il fascio $x^2+y^2-6x+k(x^2+y^2-2y)=0$.
\begin{enumerate}
\item Determina centro e raggio in funzione di $k$.
\item Trova $k$ per cui la circonferenza è tangente alla retta $x=0$.
\item Determina per quali $k$ la circonferenza ha area $4\pi$.
\end{enumerate}
\begin{solution}Centro $C\!\left(\dfrac{3}{1+k},\dfrac{k}{1+k}\right)$, $r^2=\dfrac{9+k^2}{(1+k)^2}$; tangente per $k=0$ o $k=-3$; area $4\pi$ per $k=-1\pm2$\end{solution}
\end{question}

%----------------------------------------------------------%

\begin{question}[ID=6.5, subtitle={Fascio $x^2+y^2+2x+k(x^2+y^2-4y)=0$: luogo dei centri}, difficulty=difficile, tags={fasci, circonferenze}]
Considera il fascio $x^2+y^2+2x+k(x^2+y^2-4y)=0$.
\begin{enumerate}
\item Determina il luogo dei centri al variare di $k$.
\item Trova il valore di $k$ per cui la circonferenza è tangente alla circonferenza $x^2+y^2=1$.
\item Determina per quali $k$ il centro si trova nel primo quadrante.
\end{enumerate}
\begin{solution}Centro $C\!\left(\dfrac{-1}{1+k},\dfrac{2k}{1+k}\right)$; tangente per $k=0$ o $k=-2$; $-1<k<0$\end{solution}
\end{question}
